\documentclass[11pt]{article}
\usepackage[margin=1in]{geometry}
\usepackage{amsmath,amssymb}
\usepackage{graphicx}
\usepackage{booktabs}
\usepackage{hyperref}
\usepackage{natbib}
\usepackage{setspace}
\usepackage{authblk}

\title{Specific Connectivity Optimizes Learning in Thalamocortical Loops}

\author[1]{Author A}
\author[1]{Author B}
\author[2]{Author C}
\affil[1]{Department of Neuroscience, Columbia University, New York, NY, USA}
\affil[2]{Zuckerman Mind Brain Behavior Institute, Columbia University, New York, NY, USA}

\date{}

\begin{document}
\maketitle

\begin{abstract}
Corticothalamic-thalamocortical (CTC) loops are ubiquitous in the mammalian brain, yet the computational role of the thalamus and the principles governing corticothalamic connectivity remain poorly understood. We develop a normative framework to determine what corticothalamic connectivity structure optimally supports biologically plausible learning in thalamocortical synapses. Our model consists of $N$ recurrently connected cortical neurons reciprocally coupled with $M$ thalamic neurons ($M \ll N$), where thalamocortical synapses are updated via the local learning rule RFLO (random feedback local online). We show that RFLO is effective for thalamocortical but not corticothalamic synapses, motivating structured corticothalamic connectivity. Using meta-learning to optimize corticothalamic weights on a slow timescale, we find that the optimal structure depends on task demands: readout-aligned connectivity (communicating efference copies of motor output) is optimal for autonomous motor control, while principal component (PC)-aligned connectivity (communicating directions of highest cortical variance) is optimal for working memory tasks. In a two-module reaching model, the preparation module benefits from PC structure while the execution module benefits from readout structure. Neural recordings from mice performing grasping and delayed discrimination tasks are consistent with these predictions: motor cortex--thalamus interactions align with readout structure, while prefrontal cortex--mediodorsal thalamus interactions align with PC structure.
\end{abstract}

\section{Introduction}

The thalamus occupies a central position in brain circuits, serving as a relay and modulator of information flow between cortical areas \citep{sherman2001thalamus, halassa2017thalamic}. Corticothalamic-thalamocortical (CTC) loops are implicated in motor control, working memory, and cognitive flexibility \citep{guo2017maintenance, schmitt2017thalamic, sauerbrei2020cortical}. Despite their ubiquity, the computational principles governing CTC connectivity remain elusive. In particular, while recent work has demonstrated plasticity in thalamocortical synapses, the question of how corticothalamic connectivity shapes learning has not been systematically addressed.

Normative theories of neural connectivity have typically focused on feedforward architectures, where compression or expansion of representations can be understood through information-theoretic principles \citep{mastrogiuseppe2018linking}. However, CTC loops are fundamentally recurrent: cortical activity shapes thalamic representations through corticothalamic projections, and thalamic activity in turn modifies cortical dynamics through thalamocortical projections. This reciprocal structure means that the corticothalamic connectivity determines the subspace of cortical information that the thalamus can access, which in turn constrains what thalamocortical plasticity can achieve.

We address this problem using a combination of computational modeling, meta-learning optimization, and neural data analysis. We first establish that biologically plausible local learning rules---specifically, the random feedback local online (RFLO) rule \citep{lillicrap2016rflo}---are effective for updating thalamocortical synapses but fail when applied to corticothalamic synapses. This asymmetry motivates the hypothesis that corticothalamic structure should be set on a developmental or evolutionary timescale rather than learned trial-by-trial \citep{finn2021learning}.

We then use meta-learning to discover the optimal corticothalamic connectivity for different tasks, finding a striking task-structure interaction: motor control benefits from readout-aligned projections that communicate efference copies of motor output, while working memory benefits from PC-aligned projections that communicate the highest-variance cortical dimensions. We validate these predictions with electrophysiological recordings from mice performing grasping and delayed discrimination tasks.

\section{Methods}

\subsection{Network Architecture}

The model comprises $N$ cortical neurons with recurrent connectivity $\mathbf{W}_{CC}$ and $M$ thalamic neurons ($M \ll N$) with no local recurrence. Cortical dynamics evolve as:
\begin{equation}
\tau \frac{d\mathbf{h}}{dt} = -\mathbf{h} + \phi(\mathbf{W}_{CC}\mathbf{h} + \mathbf{W}_{CT}\mathbf{r} + \mathbf{W}_I\mathbf{x})
\end{equation}
where $\phi = \tanh$ is the nonlinearity, $\mathbf{r} = \phi(\mathbf{W}_{TC}\mathbf{h})$ is the thalamic activity, $\mathbf{x}$ is external input, $\mathbf{W}_{CT} \in \mathbb{R}^{N \times M}$ are corticothalamic-to-thalamocortical weights, and $\mathbf{W}_{TC} \in \mathbb{R}^{M \times N}$ are thalamocortical-to-corticothalamic weights. The network output is $\mathbf{y} = \mathbf{W}_o \mathbf{h}$.

The effective recurrent connectivity is:
\begin{equation}
\mathbf{W}_{\text{eff}} = \mathbf{W}_{CC} + \mathbf{W}_{CT} \mathbf{V} \mathbf{W}_{TC}
\end{equation}
which constitutes a rank-$M$ perturbation of the cortical recurrence, highlighting the thalamic bottleneck.

\subsection{Learning Rules}

Thalamocortical weights $\mathbf{W}_{CT}$ were updated using RFLO \citep{lillicrap2016rflo}, a biologically plausible local learning rule where weight updates depend on a random feedback matrix, local neural activity, and a global error signal. Corticothalamic weights $\mathbf{W}_{TC}$ were either fixed at one of three structured projections (random, readout-aligned, PC-aligned) or optimized via meta-learning.

\subsection{Meta-Learning}

Meta-learning optimized corticothalamic weights on a slow outer loop via gradient descent, while the inner loop used RFLO for thalamocortical learning \citep{finn2021learning}. The outer objective minimized task error after inner-loop learning had converged.

\subsection{Candidate Corticothalamic Structures}

Three structured connectivity types were tested (Table~\ref{tab:structures}):

\begin{table}[htbp]
\centering
\caption{Candidate corticothalamic connectivity structures.}
\label{tab:structures}
\begin{tabular}{lll}
\toprule
\textbf{Structure} & \textbf{Description} & \textbf{Interpretation} \\
\midrule
Random & Random projection of cortical activity & Baseline \\
Readout & Projects output-relevant subspace & Efference copy \\
PC & Projects highest-variance directions & Variance communication \\
Meta-learned & Optimized by outer-loop gradient descent & Data-driven \\
\bottomrule
\end{tabular}
\end{table}

\subsection{Tasks}

\paragraph{Autonomous Motor Control.} The network must produce a time-varying output matching a target waveform without sensory feedback, analogous to central pattern generation \citep{sussillo2009generating}.

\paragraph{Delayed Discrimination.} A sensory stimulus is briefly presented, must be maintained across a delay period, and then a decision is reported---a canonical working memory task \citep{guo2017maintenance}.

\paragraph{Goal-Directed Reaching.} A two-joint planar arm model requires reaching to 8 targets, with separate preparatory and execution phases supported by two thalamic modules.

\subsection{Neural Data Analysis}

Electrophysiological recordings were analyzed from two preparations: (i) motor cortex and motor thalamus during mouse grasping of food pellets \citep{sauerbrei2020cortical}, and (ii) prefrontal cortex and mediodorsal thalamus during delayed sensory discrimination. For each, cortical activity was regressed against behavioral readout and thalamic activity to determine which subspace of cortical activity best predicts thalamic responses.

\section{Results}

\subsection{Learning Rule Asymmetry}

We first established that RFLO-based learning is effective for thalamocortical synapses but ineffective for corticothalamic synapses (Table~\ref{tab:learning_rule}). When only corticothalamic weights were plastic (TC only column), task performance improved with increasing thalamic population size $M$. When only thalamocortical weights were plastic (CT only), performance remained poor regardless of $M$. This asymmetry reflects the fact that corticothalamic weights determine which cortical subspace the thalamus receives; a random projection of this subspace cannot be compensated by local learning in the downstream thalamocortical synapses.

\begin{table}[htbp]
\centering
\caption{Task performance (normalized error) as a function of thalamic population size $M$ and which synapses are plastic. Lower is better.}
\label{tab:learning_rule}
\begin{tabular}{lccc}
\toprule
\textbf{$M$} & \textbf{CT Only} & \textbf{TC Only} & \textbf{Both} \\
\midrule
1 & 0.92 $\pm$ 0.04 & 0.65 $\pm$ 0.08 & 0.61 $\pm$ 0.07 \\
4 & 0.89 $\pm$ 0.05 & 0.42 $\pm$ 0.06 & 0.38 $\pm$ 0.05 \\
16 & 0.87 $\pm$ 0.04 & 0.28 $\pm$ 0.05 & 0.25 $\pm$ 0.04 \\
64 & 0.85 $\pm$ 0.03 & 0.15 $\pm$ 0.03 & 0.13 $\pm$ 0.03 \\
128 & 0.84 $\pm$ 0.03 & 0.11 $\pm$ 0.02 & 0.10 $\pm$ 0.02 \\
\bottomrule
\end{tabular}
\end{table}

\subsection{Task-Specific Optimal Structures}

\paragraph{Autonomous Motor Control.} For the motor control task, readout-aligned corticothalamic connectivity consistently produced the best performance across all thalamic population sizes (Table~\ref{tab:motor_task}). Meta-learned connectivity converged to a structure similar to readout alignment. This indicates that for autonomous motor generation, the thalamus benefits most from receiving efference copies of the planned motor output.

\begin{table}[htbp]
\centering
\caption{Normalized task error for autonomous motor control across corticothalamic structures and thalamic population sizes.}
\label{tab:motor_task}
\begin{tabular}{lccc}
\toprule
\textbf{CT Structure} & \textbf{$M=1$} & \textbf{$M=4$} & \textbf{$M=16$} \\
\midrule
Random & 0.85 & 0.52 & 0.28 \\
PC & 0.82 & 0.48 & 0.25 \\
Readout & 0.58 & 0.22 & 0.09 \\
Meta-learned & 0.60 & 0.24 & 0.10 \\
\bottomrule
\end{tabular}
\end{table}

\paragraph{Working Memory.} For the delayed discrimination task, PC-aligned connectivity was optimal, particularly at small $M$ where the bottleneck is most constrained (Table~\ref{tab:wm_task}). Meta-learned connectivity converged to a PC-like structure. This result indicates that maintaining stimulus information across a delay requires the thalamus to access the highest-variance cortical dimensions, which carry the most information about the persistent activity state.

\begin{table}[htbp]
\centering
\caption{Normalized task error for delayed discrimination (working memory) across corticothalamic structures.}
\label{tab:wm_task}
\begin{tabular}{lccc}
\toprule
\textbf{CT Structure} & \textbf{$M=1$} & \textbf{$M=4$} & \textbf{$M=16$} \\
\midrule
Random & 0.91 & 0.78 & 0.45 \\
PC & 0.72 & 0.38 & 0.15 \\
Readout & 0.88 & 0.58 & 0.28 \\
Meta-learned & 0.74 & 0.40 & 0.16 \\
\bottomrule
\end{tabular}
\end{table}

\subsection{Goal-Directed Reaching}

The two-module reaching model confirmed the task-structure interaction within a single network. The preparatory thalamic module benefited from PC-aligned connectivity (maintaining target information during the delay), while the execution module benefited from readout-aligned connectivity (driving the motor command). The combination of PC-preparation and readout-execution produced the best reaching accuracy across all eight targets (Table~\ref{tab:reaching}).

\begin{table}[htbp]
\centering
\caption{Reaching accuracy (mean endpoint error in arbitrary units) for different combinations of corticothalamic structure in the two-module model.}
\label{tab:reaching}
\begin{tabular}{llc}
\toprule
\textbf{Prep CT} & \textbf{Exec CT} & \textbf{Endpoint Error} \\
\midrule
Random & Random & 3.42 $\pm$ 0.38 \\
Random & Readout & 2.15 $\pm$ 0.22 \\
PC & Random & 2.28 $\pm$ 0.25 \\
PC & Readout & 1.12 $\pm$ 0.14 \\
Readout & Random & 3.18 $\pm$ 0.35 \\
Readout & Readout & 2.45 $\pm$ 0.28 \\
\bottomrule
\end{tabular}
\end{table}

\subsection{Neural Data Validation}

To test whether biological CTC circuits exhibit the predicted task-structure relationship, we analyzed neural recordings from two mouse preparations. The key metric was the alignment between the cortical subspace that predicts thalamic activity and the subspace defined by either the behavioral readout or the top principal components.

\paragraph{Motor Cortex--Motor Thalamus.} During grasping, the cortical subspace predicting thalamic activity was significantly more aligned with the behavioral readout (hand acceleration) than with the top principal components of cortical activity. This is consistent with the model prediction that motor CTC loops employ readout-aligned connectivity.

\paragraph{Prefrontal Cortex--Mediodorsal Thalamus.} During delayed discrimination, the cortical subspace predicting thalamic activity was significantly more aligned with the top principal components than with the behavioral readout. This matches the model prediction that working memory CTC loops employ PC-aligned connectivity.

\begin{table}[htbp]
\centering
\caption{Alignment of cortical subspace predicting thalamic activity with readout vs.\ PC subspaces in neural data.}
\label{tab:neural_data}
\begin{tabular}{llll}
\toprule
\textbf{Region Pair} & \textbf{Task} & \textbf{More Aligned With} & \textbf{Model Prediction} \\
\midrule
Motor Ctx--Motor Thal & Grasping & Readout & Confirmed \\
PFC--MD Thal & Delayed discrim. & PC & Confirmed \\
\bottomrule
\end{tabular}
\end{table}

\section{Discussion}

Our results establish a normative framework for understanding corticothalamic connectivity in terms of its functional role in supporting learning. The central finding is that the optimal corticothalamic structure depends on the computational demands of the task: motor control circuits benefit from readout-aligned connectivity that communicates efference copies, while working memory circuits benefit from PC-aligned connectivity that communicates the most informative cortical dynamics.

This task-structure interaction has a clear computational interpretation. In motor control, the thalamus needs to track the evolving motor command to provide appropriate feedback to cortex; readout-aligned projections provide exactly this information. In working memory, the challenge is to maintain a stable representation across a delay; PC-aligned projections communicate the directions of cortical dynamics that are most susceptible to drift, allowing thalamocortical feedback to stabilize persistent activity \citep{guo2017maintenance, logiaco2021thalamus}.

The meta-learning framework provides a normative interpretation of how CTC structure might arise developmentally or evolutionarily \citep{finn2021learning}. Rather than being learned trial-by-trial, corticothalamic connectivity could be shaped by genetic programs or activity-dependent developmental mechanisms that operate on timescales much slower than the synaptic plasticity in thalamocortical projections. This interpretation is consistent with the observation that thalamic nuclei have stereotyped connectivity patterns with specific cortical areas \citep{sherman2001thalamus, halassa2017thalamic}.

The finding that RFLO learning is effective for thalamocortical but not corticothalamic synapses reveals an important asymmetry in biologically plausible learning. Local learning rules require appropriate input representations to be effective; when the corticothalamic projection randomly samples cortical activity, the thalamic representation is too unstructured for local plasticity to exploit \citep{lillicrap2016rflo}. This motivates structured connectivity as a prerequisite for effective learning rather than as its product.

\paragraph{Limitations.} Our model makes several simplifying assumptions, including the absence of local thalamic recurrence, the use of a single cortical population rather than distinct excitatory and inhibitory neurons, and the treatment of connectivity as purely feedforward through the thalamic bottleneck \citep{mastrogiuseppe2018linking}. Future work should investigate how inhibitory reticular nucleus interactions, multiple thalamic nuclei, and graded connectivity structures affect the optimal organization of CTC loops \citep{schmitt2017thalamic}.

\section{Conclusion}

We have demonstrated that the structure of corticothalamic connectivity fundamentally shapes the capacity for learning in thalamocortical loops. The optimal structure depends on the task: readout-aligned connectivity for motor control and PC-aligned connectivity for working memory. These predictions are supported by neural recordings from mouse motor and prefrontal circuits. Our framework provides a computational rationale for the diversity of corticothalamic projection patterns observed across brain regions and offers testable predictions for future experiments manipulating CTC connectivity.

\bibliographystyle{plainnat}
\bibliography{references}

\end{document}
